\chapter{Cadre général et analyse des besoins}

\section{Contexte et motivation}

Les touristes et résidents qui souhaitent découvrir le Maroc se heurtent fréquemment à des informations \textbf{fragmentaires} ou \textbf{peu fiables}. Les guides généralistes et les réseaux sociaux offrent une visibilité utile mais rarement une \textbf{validation terrain} structurée. \textbf{MoroccoCheck} vise à \textbf{centraliser} la découverte des sites, à \textbf{ancrer les retours utilisateurs dans une preuve de présence} (géolocalisation) et à \textbf{encadrer la qualité} par la modération.

\section{Acteurs et périmètre fonctionnel}

Le tableau~\ref{tab:acteurs} résume les acteurs et leurs principales interactions avec le système.

\begin{table}[H]
\centering
\caption{Acteurs métier et périmètre fonctionnel}
\label{tab:acteurs}
\begin{tabular}{@{}p{3.2cm}p{10.5cm}@{}}
\toprule
\textbf{Acteur} & \textbf{Fonctionnalités principales} \\
\midrule
Touriste & Inscription, connexion (email / Google), consultation des sites, carte, check-in GPS, avis, profil, demande de passage contributeur. \\
Contributeur & Périmètre élargi de contribution communautaire (hérite des usages touriste selon le modèle de rôles). \\
Professionnel & Création ou mise à jour de fiches, revendication de site, espace professionnel, réponse aux avis. \\
Administrateur & Tableaux de bord, modération des sites et des avis, gestion des utilisateurs, traitement des demandes contributeur. \\
\bottomrule
\end{tabular}
\end{table}

\section{Besoins fonctionnels}

Les besoins suivants structurent l'application :

\begin{itemize}[leftmargin=*]
  \item \textbf{Authentification et comptes :} inscription, connexion, rafraîchissement de jetons, profil, sessions persistées.
  \item \textbf{Catalogue :} catégories, liste et détail de sites, médias, carte.
  \item \textbf{Check-in GPS :} envoi des coordonnées, calcul de distance (Haversine), contrôle de rayon et de précision, enregistrement du check-in.
  \item \textbf{Avis :} création, modification, modération, réponse éventuelle du professionnel.
  \item \textbf{Gamification :} points, badges, statistiques, classement.
  \item \textbf{Administration :} statistiques globales, files de modération, gestion des comptes.
\end{itemize}

\section{Besoins non fonctionnels}

\begin{itemize}[leftmargin=*]
  \item \textbf{Sécurité :} authentification par JWT, mots de passe hashés, validation des entrées, contrôle d'accès par rôle.
  \item \textbf{Performance :} indexation SQL, vues pour agrégats, limitation de débit optionnelle (Redis ou mémoire).
  \item \textbf{Disponibilité :} points de contrôle \texttt{/api/health}, journalisation des requêtes, suivi d'erreurs (Sentry).
  \item \textbf{Évolutivité :} architecture en couches, migrations SQL, API REST documentée par le code et les README.
\end{itemize}

\section{État de l'art succinct}

Le tableau~\ref{tab:concurrents} positionne le projet par rapport à des solutions connues (à affiner selon vos sources bibliographiques).

\begin{table}[H]
\centering
\caption{Comparaison qualitative avec des solutions existantes}
\label{tab:concurrents}
\small
\begin{tabular}{@{}p{3cm}p{5cm}p{5cm}@{}}
\toprule
\textbf{Solution} & \textbf{Forces} & \textbf{Limites} \\
\midrule
TripAdvisor & Large base d'avis, notoriété & Peu de preuve de présence sur site \\
Google Maps & Cartographie, avis & Générique, peu spécifique au positionnement métier MoroccoCheck \\
Guides / blogs & Contenu éditorial & Actualisation irrégulière, pas de validation GPS intégrée \\
\bottomrule
\end{tabular}
\end{table}

MoroccoCheck se distingue par l'association explicite \textbf{avis / présence géolocalisée} et par un \textbf{dispositif de modération} et de \textbf{rôles métiers} intégré au même produit (mobile + API + admin).
